\Askhsh{38832}{4}
\begin{ylh}{2}
    \item 3.1. Εξισώσεις 1ου Βαθμού
    \item 6.3. Η Συνάρτηση $f(x)=ax+\beta$
\end{ylh}
Στο αεροδρόμιο της πόλης έχει εγκατασταθεί ένας σταθμός με στατικά ποδήλατα γυμναστικής. Οι επιβάτες μπορούν να χρησιμοποιούν τα ποδήλατα για να φορτίζουν τις ηλεκτρικές τους συσκευές (κινητά τηλέφωνα, tablet, κλπ), παράγοντας ηλεκτρική ενέργεια με τη δική τους δύναμη. Κάθε ποδήλατο μπορεί να παράγει ενέργεια που μετατρέπεται σε ηλεκτρισμό με τη χρήση γεννήτριας. Ο σταθμός αποτελεί μέρος του προγράμματος βιωσιμότητας του αεροδρομίου, που έχει ως στόχο να ενισχύσει τη χρήση ανανεώσιμων πηγών ενέργειας και να μειώσει την κατανάλωση από το δίκτυο.

Κάθε ποδήλατο παράγει περίπου $\beta\ \si{Wh}$ (βατώρες) για κάθε $20$ λεπτά συνεχούς χρήσης. Το $80\%$ της παραγόμενης ενέργειας μετατρέπεται σε ηλεκτρική ενέργεια, χρήσιμη για φόρτιση των συσκευών, ενώ το υπόλοιπο χάνεται ως θερμότητα. Μια τυπική ηλεκτρική συσκευή χρειάζεται περίπου $10\si{Wh}$ για μία πλήρη φόρτιση.

\begin{alist}
\item Αν ένα ποδήλατο χρησιμοποιείται για $20$ λεπτά συνεχώς, πόση ενέργεια φτάνει τελικά στη συσκευή για φόρτιση;
\monades{2}
\item

\begin{rlist}
\item Αν υποθέσουμε ότι η ενέργεια που φτάνει τελικά στις συσκευές είναι ανάλογη με το πόσα εικοσάλεπτα ποδηλατούν οι επιβάτες, ποια σχέση εκφράζει την ενέργεια $(E)$ σε $\si{Wh}$ που φτάνει στις συσκευές με το πλήθος $(n)$ των εικοσαλέπτων;
\monades{5}
\item Να επιχειρηματολογήσετε υπέρ της δήλωσης: «η σχέση που συνδέει τον συνολικό αριθμό $(D)$ των συσκευών που μπορούν να φορτιστούν με την ενέργεια $(E)$, η οποία φτάνει σε αυτές, είναι της μορφής $D = a E$». Να προσδιορίσετε τη σχέση αυτή.
\monades{8}
\end{rlist}
\item

\begin{rlist}
\item Η επιβατική κίνηση του αεροδρομίου για το $2024$ ήταν $31.854.761$ επιβάτες. Αν οι συσκευές που φορτίζονται κάθε μήνα είναι $D = 780.000$ και τα εικοσάλεπτα που ποδηλάτησαν επιβάτες είναι $n = 195.000$, πόσες $\si{Wh}$ παρήγε κάθε ποδήλατο για $20\ $λεπτά συνεχούς χρήσης; Να εξηγήσετε την απάντησή σας.
\monades{8}
\item Το αεροδρόμιο πληρώνει στο δίκτυο ηλεκτρισμού $0,20$ ευρώ ανά κιλοβατώρα $(\si{kWh})$. Πόσα χρήματα εξοικονομεί κάθε μήνα χρησιμοποιώντας τον σταθμό με τα στατικά ποδήλατα;
\monades{2}

Δίνεται ότι $1\si{kWh} = 1000\si{Wh}$
\end{rlist}
\end{alist}
{\footnotesize Το παραπάνω θέμα αναπτύχθηκε στο πλαίσιο του έργου: «Ανάπτυξη Δοκιμασιών Αξιολόγησης Δεξιοτήτων Εγγραμματισμού στα μαθήματα της Νεοελληνικής Γλώσσας και Λογοτεχνίας, της Άλγεβρας, της Φυσικής και της Χημείας Α' Λυκείου Γενικού Λυκείου» Ανάδοχος: «Ειδικός Λογαριασμός Κονδυλίων Έρευνας (Ε.Λ.Κ.Ε) Πανεπιστημίου Ιωαννίνων» ΑΔΑΜ: 25SYMV016348911 2025-02-20.}

